% =====================================================================
\section{Introduction and research motivation}
\label{sec:intro}
% =====================================================================

The construction of quantum error-correcting codes from classical codes over
rings is by now a mature enterprise, and one of its most productive branches
runs through the \emph{hull} of a code: for a code $C$ equipped with a pairing
$\langle\cdot,\cdot\rangle$, the hull $C\cap C^{\perp}$ measures how far $C$ is
from being self-orthogonal, and it is precisely the hull dimension that
determines the parameters of the quantum code one can extract. Euclidean hulls
give CSS codes, Hermitian hulls give $q$-ary stabilizer codes, and --- when the
hull is small rather than large --- the entanglement-assisted formalism converts
\emph{any} classical code into a quantum code at the price of shared
entanglement.

Parallel to this, the Galois (or $\ell$-Galois) hull has been introduced over
finite fields \cite{FanZhang2017GaloisSelfDual,LiuPan2020GaloisHulls} and used
to produce entanglement-assisted codes with good parameters
\cite{Cao2021MDSGaloisHulls}. The Galois parameter $\ell$ interpolates between
the Euclidean and the Hermitian inner product, and because it can be tuned it
gives a strictly larger supply of codes with a prescribed hull dimension.

What has been missing is the passage to a \emph{non-commutative} alphabet. There
is a well-developed theory of cyclic codes over matrix rings
\cite{Alahmadi2013CyclicM2F2,LuoParampalli2018M2F2uF2,Dinh2024CyclicRingsOfMatrices,
JoshiPatelPrakash2026M4}, and there is a well-developed theory of quantum codes
from Galois hulls over commutative rings, culminating in the recent
constacyclic construction over a finite non-chain ring
\cite{ZhangKongZheng2026QuantumGaloisHulls}. The two have not met, and they do
not meet trivially: over a commutative ring the Galois twist is an
automorphism applied to one argument of the inner product, whereas over a matrix
ring the twist must be composed with the transpose to produce an
\emph{anti}-automorphism, and only then does the resulting pairing give a
perfect duality on \emph{right} modules. This paper builds that bridge and
constructs quantum codes at the far end of it.

\subsection{The chain of constructions}

Throughout, $q=p^r$, and
\[
  A_t=\F_q[u]/\langle u^t\rangle,\qquad
  R=M_2(A_t),\qquad
  R_n=R[X]/\langle X^n-1\rangle,\qquad
  \Nlen=4tn .
\]
A \emph{cyclic code of length $n$ over $R$} is a right ideal of $R_n$; this is
the right-module convention fixed for this alphabet in
\cite{Alahmadi2013CyclicM2F2}, and Lemma~\ref{lem:failure} below shows that it
is not a cosmetic choice. Under the Morita identification $R_n\cong M_2(B)$ with
$B=A_t[X]/\langle X^n-1\rangle$ \emph{commutative}, right ideals of $R_n$
correspond to $B$-submodules of $B^2$.

The contribution of this paper is the following chain, each link of which is
proved here:
\[
  \begin{gathered}
    \boxed{\text{noncommutative matrix ring}}
    \;\longrightarrow\;
    \boxed{\text{cyclic codes}}
    \;\longrightarrow\;
    \boxed{\text{Galois duality and hulls}}\\[3pt]
    \longrightarrow\;
    \boxed{\text{Gray map}}
    \;\longrightarrow\;
    \boxed{\text{quantum codes}} .
  \end{gathered}
\]
In more detail:

\begin{enumerate}[label=\textnormal{(\arabic*)},leftmargin=*]
\item \textbf{Galois duality on a non-commutative ring.} With
  $\sigma_s(a)=a^{p^s}$ the coefficient Frobenius (fixing $u$ and $X$) and
  $\theta_s=\sigma_s\circ(\cdot)^{\mathsf T}$, the pairing
  $\bb_s(x,y)=\Lambda\bigl(\sum_i\theta_s(x_i)y_i\bigr)$, where
  $\Lambda=\lambda\circ\tr$ and $\lambda$ extracts the coefficient of
  $u^{t-1}$, is non-degenerate on right $R_n$-modules, and
  $\dim C+\dim C^{\perp_s}=\Nlen$ for every code (Theorem~\ref{thm:duality}).
  We give the dual explicitly for \emph{arbitrary} finitely generated right
  ideals (Theorem~\ref{thm:dualformula}), including the reversal operator,
  which is load-bearing exactly when $n\ge3$.

\item \textbf{Hulls, with no principality hypothesis.} We prove the structural
  properties of $\Hull_s(C)=C\cap C^{\perp_s}$ for arbitrary right ideals, and
  we \emph{refute} the natural conjecture that a Galois hull is principal: hulls
  with minimal generator number $\MU_B=3$ exist
  (Theorem~\ref{thm:nonprincipal}). The phenomenon is confined:
  $\MU_B\ge3$ forces $t\ge2$ and $p\mid n$.

\item \textbf{The exact induced form.} The coefficient expansion
  $\varphi_n:R^n\to\F_q^{\Nlen}$ is not an isometry in any interesting sense,
  but it \emph{is} a form-isometry for exactly one form, and we identify that
  form: $\langle u,v\rangle^{\widehat{}}_s=u^{[s]}Jv^{\mathsf T}$ with
  $J=I_{4n}\otimes H_t$ and $H_t$ the Hankel reversal matrix
  (Theorem~\ref{thm:gram}). Duality and hulls are transported exactly, and the
  transport is valid on the non-principal hulls as well
  (Theorem~\ref{thm:grayhull}).

\item \textbf{Every Galois hull is self-orthogonal} for the transported form,
  by a one-sided argument that uses neither reflexivity nor principality
  (Theorem~\ref{thm:hullso}). This is the hinge of the whole construction: it
  is what makes the quantum step available in the general Galois regime rather
  than only in the reflexive one.

\item \textbf{Quantum codes of length $\Nlen=4tn$.} Three routes are available
  and all three are proved from scratch in symplectic language: a Galois-native
  twisted stabilizer route R1, a Euclidean/CSS route R2 and an
  entanglement-assisted route R3 giving
  $[[\Nlen,\Nlen-k-h_E,d;k-h_E]]_q$ (Theorem~\ref{thm:parameters}).

\item \textbf{A coset-distance theorem and a new bound.} The distance is the
  minimum weight of the \emph{quotient} code
  $D^{\perp_E}/\Hull_E(D)$, of dimension $K$ and length $\Nlen$
  (Theorem~\ref{thm:cosetdistance}); puncturing $\Nlen-d+1$ coordinates is
  injective on that quotient, which gives the \emph{coset Singleton bound}
  $K+d\le \Nlen+1$ (Theorem~\ref{thm:cosetsingleton}). The latter is neither
  the entanglement-assisted Singleton bound nor implied by it.
\end{enumerate}

\subsection{What the paper does and does not claim}

The codes constructed here are new in \emph{origin}, not in parameters.
Exhaustive computation over the accessible family list ($q\in\{4,8,9\}$,
$tn\le8$) certifies $28$ codes, the exact distance of each by an applicable
exact method cross-checked by an independent ambient quantum search, among them the
first stabilizer distance $4$ produced by this framework, $[[24,8,4]]_9$ and
$[[20,4,4]]_4$, and entanglement-assisted codes up to $[[20,2,10;18]]_4$. Every
stabilizer code we construct whose parameters can be compared with
\cite{GrasslCodetables} is at best equal to the best known distance --- and then
only with $d=2$ --- and every code with $d\ge3$ is one or two below it. We claim
no record. What we do claim is that the chain above is proved in the
non-commutative setting, that it produces quantum codes from hulls that are
\emph{not} principal, and that the two new bounds of item (6) hold.

\subsection{The role of the matrix-ring formulation, and what Morita
            equivalence does and does not settle}
\label{subsec:morita}

Since $X$ is central, $\Rn$ is canonically isomorphic to $M_2(\Bq)$ with
$\Bq=A_t[X]/\langle X^n-1\rangle$ \emph{commutative}, and under the Morita
correspondence right ideals of $\Rn$ correspond to $\Bq$-submodules of
$\Bq^2$; the hull of a code is then the ideal corresponding to the
intersection of the two associated submodules
(Proposition~\ref{prop:moritahull}). It is therefore fair to ask what the
matrix-ring formulation contributes. Three points.

\textbf{First, Morita equivalence is an equivalence of module categories, not
of rings with involution, and it does not carry the duality.} The duality used
throughout is defined by the anti-automorphism
$\theta_s=\sigma_s\circ(\cdot)^{\mathsf T}$ of $R$, and it is the transpose
that makes the dual of a \emph{right} ideal a right ideal with the correct
size (Lemma~\ref{lem:failure} and Theorem~\ref{thm:duality}). Over a
commutative coefficient ring the twist degenerates to $\sigma_s$ and there is
nothing to prove. The pairing of Definition~\ref{def:pairing}, the distinction
between the general Galois and the reflexive regime
(Definition~\ref{def:regimes}), and the self-orthogonality of hulls without
reflexivity (Theorem~\ref{thm:hullso}) are all statements about $R$ and
$\theta_s$ made \emph{before} any Morita translation.

\textbf{Second, what Morita does give is exactly what we use it for:} it
converts the structure theory of hulls into a computation with submodules of
$\Bq^2$. That is how the principality criterion of
Proposition~\ref{prop:moritahull} and the confinement of
Theorem~\ref{thm:nonprincipal} are proved, and how the witnesses are
constructed. \emph{We do not claim novelty from the non-commutativity of
$M_2(A_t)$ alone.} The novelty claimed is in the duality that the
non-commutative alphabet forces, and in the hull phenomena that the
$(\,t\ge2$ and $p\mid n\,)$ criterion isolates.

\textbf{Third, after the translation the residual problem is a problem about
$\Bq=\F_q[u,X]/\langle u^t,\,X^n-1\rangle$.} Its Chinese-remainder
decomposition is governed by the factorisation of $X^n-1$ over $\F_q$, and the
condition $p\mid n$ is precisely the repeated-root case. This is where the
finite-field content of the paper lives, and it is why the framework is
reported here rather than only in the language of matrices over a chain ring.

\subsection{Contributions, in order of emphasis}
\label{subsec:contributions}

The contributions, in the order in which we ask them to be read, are:

\begin{enumerate}
\item \textbf{An algebraic Galois-duality and hull framework over a
  non-commutative alphabet.} The Galois twist must be the anti-automorphism
  $\theta_s$, duality and hull theory are developed for arbitrary right ideals
  with no principality hypothesis, and self-orthogonality holds with no
  reflexivity hypothesis (Section~\ref{sec:cyclic}, Section~\ref{sec:hulls}).
\item \textbf{Non-principal hull structure, and a counterexample.} Galois
  hulls need not be principal; $\MU_{\Bq}\ge3$ forces $t\ge2$ and $p\mid n$;
  and no closed formula of gcd type exists
  (Theorem~\ref{thm:nonprincipal}, Theorem~\ref{thm:noclosed}).
\item \textbf{A duality-preserving expansion to $\F_q^{\Nlen}$.} The
  transported form is identified exactly, with Hankel Gram matrix
  $I_{4n}\otimes H_t$, and duality and hulls are preserved, non-principal ones
  included (Section~\ref{sec:gray}).
\item \textbf{Quantum stabilizer and entanglement-assisted constructions.}
  Three routes of length $\Nlen=4tn$, with the characteristic-two obstruction
  of Theorem~\ref{thm:obstruction} delimiting the Euclidean route
  (Section~\ref{sec:quantum}).
\item \textbf{Exact distance certificates and two bounds.} The coset-distance
  theorem and the coset Singleton bound $K+d\le\Nlen+1$, and the certification
  of $28$ codes, each by an exact method cross-checked independently
  (Section~\ref{sec:cert}).
\end{enumerate}

The enumerators and dimension-distribution statistics produced by the scan are
\textbf{supporting material}. They are reported where they support a
construction and are not presented as principal results.

\subsection{On the parameters obtained}
\label{subsec:parametersposition}

Every stabilizer code constructed here is at best \emph{equal} to the best
known distance, and then only at $d=2$; every stabilizer code with $d\ge3$ is
one or two below the best known value (Section~\ref{sec:cmp}). The
entanglement-assisted codes carry the largest distances obtained in this
framework, and they cannot be benchmarked at all, because no
entanglement-assisted table is available; the codes over $\F_8$ are likewise
unbenchmarked. \textbf{No record parameter is claimed.} What is offered is an
exact algebraic \emph{origin} for these codes: non-commutative Galois hulls,
transported to $\F_q$ with duality and hulls preserved, with non-principal
hulls admitted, and with two new bounds constraining the parameters that
result. The value of the construction is that every parameter reported is
derived and certified rather than merely asserted, and that the algebraic
hypotheses behind each tuple are visible.


\subsection{Organisation}

Section~\ref{sec:lit} reviews the literature and states the research gap.
Section~\ref{sec:prelim} fixes the rings, the right-module convention and the
anti-automorphism, and proves the failure lemma that motivates the convention.
Section~\ref{sec:cyclic} develops cyclic codes and Galois duality.
Section~\ref{sec:hulls} is the algebraic heart: hull structure, the
non-principal counterexample and its confinement.
Section~\ref{sec:gray} transports everything to $\F_q^{\Nlen}$ and identifies
the induced form. Section~\ref{sec:quantum} builds the quantum codes.
Section~\ref{sec:cert} proves the coset-distance theorem and the coset Singleton
bound, describes the certification methodology and reports the parameter
analysis. Section~\ref{sec:cmp} compares with the literature and states the
novelty precisely. Section~\ref{sec:conclusion} collects open problems. Appendix~A
contains the full parameter tables and Appendix~B the reproducibility package.
